\section*{第 \ref{ch:g9:our-responsibility} 章——我们的责任}

\begin{solution}{exo:g9:our-responsibility:1}
将日常作息视为对器官的合理维护；
为有害物质设卡，并拒之门外；
利用生殖图谱及其中的地址；并把
家族倾向读作一份维护日程表。
\end{solution}

\begin{solution}{exo:g9:our-responsibility:2}
传播（\cref{ch:g9:microbes-and-infection}）——我的
\omterm{def:g1:healthy-habits:hygiene}{卫生}保护你；群体\omterm{def:g9:immune-defenses:memory}{免疫}
（\cref{ch:g9:immune-defenses}）——我的接种保护
无法接受\omterm{def:g9:immune-defenses:memory}{免疫}训练者；医疗公共资源
（\cref{ch:g9:how-species-change}中的抗性、
\cref{ch:g9:unique-individuals}中的\omterm{prop:g4:journey-of-food:absorb}{血液}登记）——
靠个人选择维持运转，或得以补充储备。
\end{solution}

\begin{solution}{exo:g9:our-responsibility:3}
因为\omterm{def:g9:heredity-and-traits:heredity}{遗传}倾向是生活其中的范围，而非注定的命运：
\omterm{def:g7:heart-and-circulation:vessels}{动脉}易硬化的家族，若能据此调整饮食、活动和
检查，便能改写结局——知识把这笔交易
变为一份习惯与
体检的日程表。
\end{solution}

\begin{solution}{exo:g9:our-responsibility:4}
例如：\omterm{prop:g4:heart-and-blood:pulse}{脉搏}与血压——在路边即可读出的循环状况
（\cref{ch:g7:heart-and-circulation}）；叩诊锤——
受检的\omterm{def:g8:nervous-communication:neuron}{神经}回路
（\cref{ch:g8:nervous-communication}）；接种记录——
储备的记忆驻军
（\cref{ch:g9:immune-defenses}）。
\end{solution}

\begin{solution}{exo:g9:our-responsibility:5}
\omterm{def:g5:food-webs:web}{食物网}的缓冲作用依赖众多\omterm{def:g5:biodiversity:species}{物种}——
网中的节点与替补者；
而选择的未来依赖于多样性——
是每一次\omterm{def:g6:exploring-our-environment:components}{环境}变化都可取用的储备。
二者都是可以动用、积存的功能，而非装饰。
\end{solution}

\begin{solution}{exo:g9:our-responsibility:6}
先观察，后下结论；
以一个差异加一组对照进行公平的检验；
在采信某种说法前，先说明其机制；
权衡相互独立的证据；尊重体内外的生命机器。
\end{solution}

\begin{solution}{exo:g9:our-responsibility:7}
发热的学生留在家中，便从源头切断
疫情的传播——停课一天就是切断的途径，
比日后任何药物都更省。
（\cref{pb:g9:microbes-and-infection:1}中的箴言
在实践中的体现）。
\end{solution}

\begin{solution}{exo:g9:our-responsibility:8}
群体的防线在于：每个接受过训练的人
都是切断传播的死胡同。被依赖者包括：过小的婴儿
（\cref{pb:g9:immune-defenses:1}中受看护的接触者）
以及\omterm{def:g9:immune-defenses:memory}{免疫}力低下的患者，他们只能依靠
周围人的\omterm{def:g9:immune-defenses:memory}{免疫}训练来获得保护。
\end{solution}

\begin{solution}{exo:g9:our-responsibility:9}
“每一种\omterm{prop:g3:sorting-living-things:groups}{昆虫}”包括果园\omterm{prop:g3:flowers-fruits-seeds:transform}{果实}所依赖的\omterm{def:g5:flower-to-fruit:steps}{传粉}者，
以及替人除害却不计报酬的食虫\omterm{prop:g3:sorting-living-things:groups}{昆虫}
（\cref{ex:g5:biodiversity:orchard}）；
审计预测，替补者死亡之后，人工授粉将要付费，
病虫害也会更严重——全面的“保护”若清除了这支
劳动力，不过是把账单推迟，而非真正的保护。
\end{solution}

\begin{solution}{exo:g9:our-responsibility:10}
每一次断续、随意的用药，都在病区运行一次选择的筛子——
杀死敏感者，促进抗性者繁殖，
因此每一次滥用，都在消耗全体人的药物未来效用。
应养成的习惯：按规定完成疗程，
且不为\omterm{def:g9:microbes-and-infection:sorted}{病毒}性疾病要求使用抗生素。
\end{solution}

\begin{solution}{exo:g9:our-responsibility:11}
观察：这篇帖子究竟展示了什么——是证词与光鲜的说法，还是数据？
检验：是否有过公平的试验——一个差异、一组对照，
还是仅凭一个感激的病例？
机制：所声称的疗效阻断了哪一个环节——能否说清？
见证：独立来源是否一致，还是一切都溯源于卖家？
尊重：它是否要求你为了一个奇迹
而放弃行之有效的维护？
五步之中，这类帖子往往
经不起前两步的检验。
\end{solution}

\begin{solution}{exo:g9:our-responsibility:12}
例如：“每一种\omterm{def:g6:exploring-our-environment:components}{生物}都由\omterm{def:g5:first-look-at-cells:cell}{细胞}构成，
运行着同一套古老的化学机制——你也不例外。你的
身体是一台由器官、边界与信使组成的维持机器，
而它的维护正是你自己的责任。每一餐、每一次\omterm{def:g7:what-is-respiration:respiration}{呼吸}，
都把你连入由\omterm{def:g5:food-webs:roles}{生产者}、\omterm{def:g5:food-webs:roles}{消费者}与\omterm{def:g6:decomposers-and-soil:decomposers}{分解者}组成的网络——
维系这张网，同样是你的责任。整体是同一个\omterm{def:g1:plants-around-us:tree}{分支}繁茂的家族，
写在同一套密码之中，
在\omterm{prop:g9:how-species-change:selection}{自然选择}下不断变化——
而你，正是那个读懂这棵树的\omterm{def:g1:plants-around-us:tree}{分支}。”
\end{solution}

\begin{solution}{pb:g9:our-responsibility:1}
\textbf{1.} 飞沫与双手——因此要掩口咳嗽、勤洗手、
通风换气，生病时待在家中。员工应当优先，
因为他们连着这栋楼里最脆弱的人群：
他们的训练恰好覆盖了无法受训者聚集之处。

\textbf{2.} 登记册储存着配型相符的\omterm{prop:g4:journey-of-food:absorb}{血液}，以备
意外事故与手术之需；账本储存着受过训练的
记忆。抽屉：人人可取——登记册的配型按群体查找，
训练池的盾牌覆盖年轻人、老年人
与病患者。

\textbf{3.} 例如：“某病房曾仅凭洗手，
就使产妇死亡数减半：所谓的‘麻烦’打断了肉眼看不见的信使，
而自塞麦尔维斯以来的一个半世纪，
只是证实了一件事：世上最便宜的药
竟出自一只
肥皂盘。”

\textbf{4.} 抗生素这一公共资源，在参差不齐的选择中，
每个未完成的疗程都会放过并促进耐药的少数个体，
把药效白白耗在每个人未来的
\omterm{def:g9:microbes-and-infection:stages}{感染}之上。

\textbf{5.} 这种喷雾会连同害虫一并夺走\omterm{def:g5:flower-to-fruit:steps}{传粉}者与食虫者：
坐果率随蜜蜂减少而下降，
幸存害虫的下一次爆发将遇不到
瓢虫，也遇不到山雀——免费的服务停止了，
账单则以租来的蜂箱和更多的喷雾形式回来。

\textbf{6.} 豆绿色的水华在温暖而静止的水中整夜进行\omterm{def:g7:what-is-respiration:respiration}{呼吸作用}，
而这种水中的氧很少：黎明时的
溶氧最低值降至谷底
（见\cref{ex:g7:oxygen-in-the-water:curve}）。处方如下：
夜间运行的喷泉，疏浚淤泥与过度生长的水草——
同时留意曾经滋养这场水华的营养输入。

\textbf{7.} \omterm{def:g2:where-animals-live:habitat}{栖息地}：这片沼泽是繁殖地——
没有\omterm{def:g2:where-animals-live:habitat}{栖息地}，就没有居民。\omterm{def:g5:food-webs:web}{食物网}方面：这里的\omterm{def:g5:biodiversity:species}{物种}\omterm{def:g5:flower-to-fruit:steps}{传粉}、
净化水体，并滋养周边的\omterm{def:g3:food-chains:chain}{食物链}。保护区方面：它的
多样\omterm{def:g8:reproduction-and-environment:population}{种群}储存着艰难时期将会用到的
\omterm{def:g9:genes-and-dna:gene}{等位基因}。（鹳鸟回到草甸所在之处。）

\textbf{8.} 未付的账单：洪水被蓄存并缓慢释放，
水体得到过滤，\omterm{def:g5:flower-to-fruit:steps}{传粉}者得到安置，
幼鱼得到照料——这些服务只有在
沼泽被排干、城镇转向工程师购买时才被计价。

\textbf{9.} 答案不限，一组可以站得住脚的组合：最好的——
保护区（在三条可持续利用准则下资本都已存入）；
最差的——未完成的疗程，或喷雾方案
（公共资源被耗尽，工人被迫离开），依据
“不加速攫取”“留住工人”“归还所借”三条\omterm{def:g3:sorting-living-things:criterion}{标准}来判断。

\textbf{10.} 观察到：黎明时有死鱼，水体发绿，
读数偏低。公平的检验：将黎明与午后的探测进行比较，
将这个池塘与一个清澈的池塘进行比较。已指明的机制：夜间
水华的\omterm{def:g7:what-is-respiration:respiration}{呼吸作用}，处在温水下降的
溶氧上限之下。佐证：探测数据、普查中的耐受者，
幸存者与时间线，三者互相印证。可取之处在于：
处方调控的是池塘中的\omterm{def:g6:exploring-our-environment:components}{生物}机制，
而不是只对症状用药。

\textbf{11.} 例如：“我们的流感浪潮止于洗手与
病假落实之处——健康是一笔
共同的账目。我们所拯救的沼泽是资本，而非
风景，它将以蓄住的洪水与授粉的收成回报我们。
今年声音最响的那些说法，经不起这一问：
请说出机制。保留这一方法吧：它
是小镇最便宜的基础设施。”

\textbf{12.} 答案不限——一种诚实的可能是：
“蜗牛与你由同样的四个字母构建而成，
请守护整个家族。”
\end{solution}
